\documentclass[12pt]{article}
\usepackage[utf8]{inputenc}
\usepackage[T1]{fontenc}
\usepackage{hyperref}
\title{Vinekula Language Resources for K-12 Curriculum Development}
\author{Compiled by Claude Sonnet 5 (claude-sonnet-5) for Andr\'as Kornai}
\date{}
\begin{document}
\maketitle

\section{Introduction}
Vinekula (ISO 639-3 \texttt{tmi}) is also called Tutuba. It is an Oceanic Austronesian language of Vanuatu, spoken on Tutuba Island and the southeast tip of Espiritu Santo. Wikipedia gives about 500 speakers. Very little material exists.

\section{Primary Educational Resources}
None found. Endangered Languages reports that the language has no literacy tradition.

\section{Linguistic and Reference Materials}
\begin{itemize}
\item Wikipedia, Tutuba language: \url{https://en.wikipedia.org/wiki/Tutuba_language}
\item Endangered Languages Project entry: \url{https://endangeredlanguages.com/lang/2447}
\end{itemize}

\section{Religious and Community Resources}
Joshua Project reports that Scripture translation has started. I could not confirm any published text. \url{https://joshuaproject.net/people_groups/15664}

\section{Digital and Audio Resources}
None found.

\section{Curriculum Development Context}
There are no known school materials and no orthography was found in the search results. Any curriculum would need primary documentation work with the community first.
\end{document}